\appendix
\section{Artifact Appendix}

\subsection{Abstract}
The artifact provides the paper-relevant MBPriorQ software, representative
Qwen3-0.6B experiments on WikiText2, and curated MBPriorQ accelerator sources in
SpinalHDL. The public
workflows reproduce the representative W4A4 accuracy and central ablations,
validate mask/scale accounting, and functionally exercise the packet scheduler,
regular/refined MultiMSA, shared FPU pool, output synchronization, and
packetization with Verilator.
Commercial EDA data and reproductions of comparison methods are intentionally
outside the artifact scope.

\subsection{Artifact Check-list (Meta-information)}
{\small
\begin{itemize}\setlength\itemsep{0pt}\setlength\parskip{0pt}
  \item \textbf{Algorithm:} MBPriorQ W4A4 fake quantization, online activation
    prior, selective 16-to-4 refinement, and EBW accounting.
  \item \textbf{Program:} Python/PyTorch software; Scala/SpinalHDL hardware;
    Bash/Python automated validation.
  \item \textbf{Model/data:} public Qwen3-0.6B and WikiText2 raw test; the small
    importance matrix is included, but model weights and WikiText2 are not.
  \item \textbf{Run-time environment:} Linux x86-64; Python 3.12, PyTorch 2.10,
    Transformers 4.57; Java 17, sbt 1.10.2, SpinalHDL 1.8, Verilator 4.034+.
  \item \textbf{Hardware:} core PPL: NVIDIA GPU with at least 8\,GB free VRAM;
    functional/hardware simulation: CPU only. Tested on an RTX 5090.
  \item \textbf{Metrics/output:} perplexity, effective bit width, cycle counts,
    and generated protocol CSVs.
  \item \textbf{Disk/time:} about 10\,GB for the core workflow; under 1 minute
    for software smoke, tens of minutes for Qwen3-0.6B experiments, about one
    minute for hardware modules, and 8--15 minutes for the public top.
  \item \textbf{Availability:} tagged release \ArtifactVersion; DOI
    \ArtifactDOI; source license \ArtifactLicense.
\end{itemize}
}

\subsection{Description}
\subsubsection{How to Access}
The immutable artifact is distributed through the DOI above. The evaluated Git
tag and archive version must match. The repository root \path{README.md}
indexes the claim map, provenance, expected outputs, and release checksums.

\subsubsection{Hardware, Software, Data, and Models}
The functional tensor check needs a Linux host with approximately 8\,GB RAM.
The representative PPL run needs at least 16\,GB system RAM and 8\,GB free CUDA
VRAM. Open hardware simulation additionally requires Java 17, sbt, Verilator, Make,
and a C++ compiler; no GPU is required. Conda specifications are provided in
\path{environment/}. Qwen3-0.6B and WikiText2 are downloaded from their
official providers under their respective terms. The archive retains the Qwen
model-license notice but does not redistribute model weights or WikiText2.

\subsection{Installation}
From the repository root, create the software environment with:
\begin{quote}\small\ttfamily
conda env create -f environment/software.yml\\
conda activate mbpriorq-ae\\
python -m pip install -e software
\end{quote}
For RTL, follow \path{environment/README.md} to create the hardware environment
and install sbt 1.10.2. Installation takes about 10--20 minutes with a normal
package cache and Internet connection.

\subsection{Experiment Workflow}
Download \path{Qwen/Qwen3-0.6B} and set \path{MODEL_PATH} before steps 2--3.
Run the following commands from the repository root:
{\scriptsize
\begin{flushleft}\setlength\parskip{0pt}
1. \path{scripts/run_smoke.sh}\\
2. \path{experiments/core_accuracy/run.sh}\\
3a. \path{experiments/activation_attribution/run.sh}\\
3b. \path{experiments/granularity_ablation/run.sh}\\
4a. \path{scripts/run_hardware_modules.sh}\\
4b. \path{scripts/run_hardware_system.sh}
\end{flushleft}
}
Step 1 checks tensor and EBW semantics. Steps 2--3 reproduce PPL/ablation rows
and reuse generated checkpoints locally. Step 4 validates actual SpinalHDL
modules and the public 1024-bit packet top against recorded functional outputs.

\subsection{Evaluation and Expected Results}
The full core run must process 146 contiguous 2048-token windows and report:
\begin{center}
\small
\begin{tabular}{lrr}
\toprule
Method & PPL & Act. EBW \\
\midrule
BF16 & 20.9240 & -- \\
W4A4 MBPriorQ & 24.2289 & 4.994211 \\
\bottomrule
\end{tabular}
\end{center}
For 16-to-8/4/2 refinement, expected PPL is 24.8590/24.2289/23.9532 and
activation EBW is 4.704345/4.994211/5.613351. Each script enforces its recorded
tolerance and sample count. Hardware checks must reproduce all expected CSVs.

\subsection{Experiment Customization}
\texttt{NUM\_SAMPLES} shortens PPL runs for functional debugging; only zero
(all windows) is a reproduced full row. \texttt{DATASET\_PATH} selects a local
WikiText2 copy. The granularity workflow exposes 16-to-8/4/2 configurations.
Hardware scripts may be run separately to isolate module tests from the heavier
public-top Verilator compilation.

\subsection{Notes}
The complete 19-entry accuracy sweep and large-model extensions are outside the
default AE rerun; the public Results Reproduced target uses the fully executable
Qwen3-0.6B representative workflow. Hardware tests target functional validity,
not speedup, area, or power. Commercial synthesis evidence and reproductions of
comparison methods are not distributed or evaluated by this artifact.

\subsection{Methodology}
Artifact review and badging follow the ACM policy at
\url{https://www.acm.org/publications/policies/artifact-review-and-badging-current}
and the cTuning AE methodology at \url{https://ctuning.org/ae}.
